\documentclass[11pt, letterpaper]{article}

\usepackage[margin=1in]{geometry}
\usepackage{hyperref}
\usepackage{titlesec}
\usepackage{enumitem}
\usepackage{parskip}
\usepackage{amsmath}
\usepackage{amssymb}

\hypersetup{
    colorlinks=true,
    linkcolor=blue,
    urlcolor=cyan,
    pdftitle={Technical Design: VR Health Data Transmission}
}

\title{\textbf{Technical Design Report} \\ \Large Secure Bi-Directional Health Data Transmission from VR Edge Devices to Web Portals}
\author{}
\date{\today}

\begin{document}

\maketitle

\begin{abstract}
\noindent This report outlines the technical, cybersecurity, and regulatory strategies required to transmit bi-directional health data from a VR headset connected to public or private Wi-Fi to a clinical web portal. The architecture evaluates cloud-native, on-premises, and hybrid approaches, emphasizing compliance with HIPAA, 42 CFR Part 2, and modern zero-trust cybersecurity principles.
\end{abstract}

\tableofcontents
\newpage

\section{Introduction}

\subsection{Project Scope}
The objective is to establish a secure data pipeline: VR Headset (edge device) $\rightleftarrows$ Cloud Backend $\rightleftarrows$ Web Portal (clinician/patient). The target hardware is a Pico Android-based edge device, representing an AIO VR environment. The system functions as a potentially FDA-regulated Software as a Medical Device (SaMD) or a Class II device.

\subsection{Pre-existing Technologies}
The proposed data flow mirrors established healthcare IoT and wearable technologies, which securely transmit telemetry and Protected Health Information (PHI) to cloud environments:
\begin{itemize}
    \item Diabetes insulin pumps and Dexcom sensors (skin-worn IoT) leveraging BLE to phone to cloud pipelines.
    \item Dexcom Clarity cloud ecosystems.
    \item Aura rings for biometrics.
\end{itemize}

\section{System Architecture}

\subsection{Architecture Patterns}
Three primary architectural patterns were evaluated for this deployment:

\textbf{Cloud-Native (AWS/Azure)}
This is the recommended approach for scalability and alignment with modern HIPAA-compliant cloud offerings. The pipeline follows: Pico headset $\rightarrow$ Wi-Fi (WPA3-Enterprise or private network) $\rightarrow$ TLS 1.3 API Gateway $\rightarrow$ Cloud WAF $\rightarrow$ Microservices backend (containers/serverless) $\rightarrow$ Encrypted database $\rightarrow$ Web portal.

\textbf{On-Premises / Private Cloud}
While this reduces latency (critical for VR) and minimizes initial PHI exposure, it is highly costly and lacks scalability. It is generally avoided unless mandated by specific hospital infrastructure.

\textbf{Hybrid Architecture}
Session data is processed at the edge on the device. Only aggregated or de-identified telemetry is sent to the cloud in real-time, while full PHI is batch-synced over a secure tunnel at defined intervals. This provides an excellent "privacy by design" regulatory story but increases system complexity.

\subsection{Technology Stack}
\begin{itemize}
    \item \textbf{Device (Edge):} Android (Kotlin/Java), Unity (VR app), and Local SQLite (encrypted).
    \item \textbf{Suggested Protocols:} REST for simple transactions, gRPC for low-latency VR sync, and MQTT for IoT-style data streaming.
    \item \textbf{Backend (AWS):} API Gateway, Lambda/ECS for microservices, RDS (PostgreSQL) for structured data, and S3 for object storage.
    \item \textbf{Frontend:} React combined with Next.js, strictly served over HTTPS.
\end{itemize}

\section{Cybersecurity Strategy}

\subsection{Device Security (VR Headset)}
Security at the edge is paramount. The device must enforce secure boot and utilize signed firmware. A Trusted Execution Environment (TEE) should be used alongside an encrypted local SQLite database to protect data at rest on the headset. Fail-safe modes must be implemented, allowing the device to work offline and sync data securely when a connection is re-established.

\subsection{Network Security}
All data in transit must be protected by TLS 1.3. Mutual TLS (mTLS) should be established between the device and the backend to verify device identity. For highly sensitive synchronization events, dedicated VPN tunnels are recommended.

\subsection{Backend and Web Portal Security}
The infrastructure must adopt a Zero Trust Architecture. 
\begin{itemize}
    \item \textbf{Identity and Access:} Managed via OAuth 2.0 / OpenID Connect, utilizing AWS Cognito.
    \item \textbf{Authorization:} Role-Based Access Control (RBAC) must be enforced, alongside mandatory Multi-Factor Authentication (MFA) for clinician access.
    \item \textbf{Threat Mitigation:} Deployment of a Web Application Firewall (WAF), rigorous API rate limiting, and strict input validation to prevent injection attacks. Key management is handled via AWS KMS.
\end{itemize}

\section{Regulatory Compliance}

\subsection{HIPAA and General Standards}
The system must maintain compliance with HIPAA, NIST, and FIPS standards. This requires guaranteed end-to-end encryption for all PHI and Personally Identifiable Information (PII) both at rest and in transit. The architecture must enforce strict access controls and maintain immutable audit trails of all data access.

\subsection{42 CFR Part 2 (Substance Use Disorder)}
Handling Substance Use Disorder (SUD) records involves highly restrictive compliance requirements under 42 CFR Part 2. Data segmentation is critical. SUD behavioral data must be explicitly tagged, isolated, and separated from general medical data and user identifiers. This requires separate storage mechanisms or strictly delineated logical access controls within the database.

\subsection{Privacy by Design}
The system adheres to the principle of data minimization. Raw VR data should not be streamed unless strictly necessary; edge processing should be utilized to extract relevant metrics. Data must be de-identified at the earliest possible stage in the pipeline.

\end{document}
